% ==============================================================
\section{問題定義}\label{sec:problem}
% ==============================================================

本節では、まず Apriori-Window~\cite{aprioriwindow}が扱う1次元設定を
整理し（\ref{sec:1d_review}節）、
続いてそれを $d$ 次元に拡張した本研究の問題定義を与える（\ref{sec:nd_extension}節）。

% ==============================================================
\subsection{Apriori-Window（1次元）の復習}\label{sec:1d_review}
% ==============================================================

$\mathcal{I} = \{i_1, \ldots, i_m\}$ をアイテム集合、
$\mathcal{D}_{\mathrm{1D}} = \{(t_n, d_n)\}_{n=1}^N$ を
タイムスタンプ付きトランザクション列とする。
$t_n \in \mathbb{Z}_{\geq 0}$ は整数タイムスタンプ（または離散化済みビン）、
$d_n \subseteq \mathcal{I}$ はアイテム集合である。

パターン $P \subseteq \mathcal{I}$ の\emph{出現集合}を
$\mathrm{occ}(P) = \{t_n : P \subseteq d_n\}$ と定義する。

\begin{definition}[区間サポート~\cite{aprioriwindow}]\label{def:sup1d}
ウィンドウ幅 $W_t$ のもとで、ウィンドウ左端 $l$ に対する
パターン $P$ の\emph{区間サポート}を
\[
  \mathrm{sup}(P,\, l,\, W_t)
  = \bigl|\{t \in \mathrm{occ}(P) : l \leq t < l + W_t\}\bigr|
\]
と定義する。
\end{definition}

\begin{definition}[密集窓・密集区間~\cite{aprioriwindow}]\label{def:dense1d}
サポート閾値 $\theta$ に対し、
$\mathrm{sup}(P, l, W_t) \geq \theta$ を満たすとき、
窓 $[l, l+W_t]$ を $P$ の\emph{密集窓}と呼ぶ。
密集窓の開始点 $l$ が連続する最長区間 $[s, e]$ を
$P$ の\emph{密集区間}と呼ぶ。
\end{definition}

Apriori-Window~\cite{aprioriwindow}は、
$t$ 軸ストライドによる不要スキャンの排除と
候補区間の交差に基づく Apriori 枝刈りにより、
全パターンの密集区間を効率的に列挙する。

% ==============================================================
\subsection{$d$ 次元への拡張}\label{sec:nd_extension}
% ==============================================================

本研究では、各トランザクションが時刻 $c_0$（時間軸）に加えて
$d-1$ 個の追加座標 $(c_1, \ldots, c_{d-1})$（空間・カテゴリなど）を持つ
\emph{$d$ 次元トランザクションデータ}を対象とする。

\paragraph{記法.}
$\mathcal{D} = \{(\bm{c}_n, d_n)\}_{n=1}^N$ を
\emph{$d$ 次元トランザクション列}とする。
ここで $\bm{c}_n = (c_n^{(0)}, c_n^{(1)}, \ldots, c_n^{(d-1)}) \in \mathbb{R}^d$ は
座標ベクトル（第0軸は時間）、
$d_n \subseteq \mathcal{I}$ はアイテム集合である。

\begin{definition}[$d$ 次元ビンインデックス]\label{def:bin_nd}
ビン幅ベクトル $\bm{\Delta} = (\Delta_0, \ldots, \Delta_{d-1})$ のもとで、
トランザクション $(\bm{c}, d)$ のビンインデックスを
\[
  \bm{b}(\bm{c}) = \bigl(\lfloor c^{(k)} / \Delta_k \rfloor\bigr)_{k=0}^{d-1}
\]
と定義する。グリッドサイズを $G = \prod_{k=0}^{d-1} N_k$ とする
（$N_k$ は第 $k$ 次元のビン数）。
\end{definition}

\begin{definition}[$d$ 次元アイテム出現グリッド]\label{def:item_grid_nd}
パターン $P \subseteq \mathcal{I}$ の\emph{出現グリッド} $I_P$ を
\[
  I_P[\bm{b}]
  = \bigl|\{n : \bm{b}(\bm{c}_n) = \bm{b},\; P \subseteq d_n\}\bigr|
\]
と定義する。すなわち、ビン $\bm{b}$ において
パターン $P$ を含むトランザクション数である。
\end{definition}

\begin{definition}[$d$ 次元スライディングウィンドウサポート]\label{def:support_nd}
ウィンドウビン幅ベクトル $\bm{W} = (W_0, W_1, \ldots, W_{d-1})$ のもとで、
グリッド点 $\bm{s} = (s_0, s_1, \ldots, s_{d-1})$ を開始点とする
ウィンドウ内のパターン $P$ の\emph{サポート}を
\begin{equation}\label{eq:support_nd}
  S_P(\bm{s})
  = \sum_{\bm{b} \in \prod_{k=0}^{d-1} [s_k,\, s_k + W_k)}
    I_P[\bm{b}]
\end{equation}
と定義する。全グリッド点 $\bm{s}$ にわたる $S_P(\bm{s})$ を
$P$ の\emph{サポートサーフェス}と呼ぶ。
\end{definition}

\begin{remark}[1次元設定への帰着]\label{rem:1d_reduction}
$d = 1$（時間軸のみ）のとき、$S_P(s_0) = \mathrm{sup}(P, s_0, W_0)$
（定義~\ref{def:sup1d}）に帰着する。
すなわち本研究の問題設定は Apriori-Window~\cite{aprioriwindow}の真の拡張である。
さらに $d = 3$（時間 $\times$ 空間 $x \times$ 空間 $y$）のとき、
本稿で示す3次元実験結果が得られる。
\end{remark}

\begin{definition}[密集セル]\label{def:dense_cell_nd}
サポート閾値 $\theta$ に対し、
グリッド点 $\bm{s}$ が
$S_P(\bm{s}) \geq \theta$ を満たすとき、
$\bm{s}$ を\emph{密集セル}と呼ぶ。
パターン $P$ の密集セル集合を
$\mathcal{L}_\theta(S_P)$ と表す。
\end{definition}

\begin{definition}[密集領域]\label{def:dense_region_nd}
密集セル集合 $\mathcal{L}_\theta(S_P)$ の
$2d$-近傍連結成分（各次元の正負方向 $\pm \bm{e}_k$ を隣接と定義）の
各成分を\emph{密集領域}と呼ぶ。
$d = 3$ のとき $2d = 6$ 近傍（辺共有隣接）となる。
\end{definition}

1次元の「密集区間」（定義~\ref{def:dense1d}）は、
$d = 1$ のときの密集領域に対応する。

\paragraph{Apriori 包含定理.}

\begin{lemma}[密集領域の単調性]\label{lem:apriori}
任意のパターン $P$ とその部分集合 $Q \subset P$ に対して、
$\mathcal{L}_\theta(S_P) \subseteq \mathcal{L}_\theta(S_Q)$ が成立する。
\end{lemma}

\begin{proof}
任意のビン $\bm{b}$ に対し、$P \subseteq d_n$ ならば $Q \subseteq d_n$ であるから
$I_P[\bm{b}] \leq I_Q[\bm{b}]$。
接頭辞和の単調性より $S_P(\bm{s}) \leq S_Q(\bm{s})$。
\end{proof}

補題~\ref{lem:apriori}は Apriori-Window の候補区間交差枝刈り~\cite{aprioriwindow}の
$d$ 次元版であり、$k \geq 2$ の枝刈りの根拠である（\ref{sec:alg_k2}節）。

\subsection{$d$ 次元密集領域検出問題}

\begin{problem}[$d$ 次元密集領域検出]\label{prob:dense_detection}
\textbf{入力}：$d$ 次元トランザクション列 $\mathcal{D}$、
ビン幅 $\bm{\Delta} \in \mathbb{R}^d$、
ウィンドウビン数 $\bm{W} \in \mathbb{Z}_{>0}^d$、
サポート閾値 $\theta$、最大パターン長 $L_{\max}$。

\textbf{出力}：
各パターン $P$ に対し、$P$ の密集領域の全集合
$\{\mathcal{R}_1, \mathcal{R}_2, \ldots\}$（各 $\mathcal{R}_j$ は密集セルの連結成分）。
\end{problem}

問題~\ref{prob:dense_detection}の計算量は
サポートサーフェス $S_P$ の計算コストに支配される。
単純実装の $O\!\left(M \cdot G \cdot \prod_{k=0}^{d-1} W_k\right)$ を
提案アルゴリズム（\ref{sec:method}節）では
$O\!\left(M \cdot G + N \cdot \prod_{k=1}^{d-1} W_k\right)$ に削減する
（第0軸を時間軸 $t$ とする）。
